\documentclass[10pt,a4paper]{amsart}
\usepackage[margin=19mm]{geometry}
\usepackage{amsmath,amssymb,mathtools}
\usepackage[scr=rsfs,cal=euler]{mathalfa}
\usepackage{xfrac}
\usepackage[unicode,hidelinks,bookmarks=false]{hyperref}
\newcommand{\ie}{i.e.\ }
\newcommand{\cf}{cf.\ }
\newcommand{\resp}{respectively\ }
\newcommand{\Subsection}{\subsection}
\providecommand{\nobreakdash}{\nobreak\hbox{-}}
\setlength{\emergencystretch}{2em}
\multlinegap0pt
\usepackage{bm}
\usepackage{bbm}
\usepackage{dsfont}
\usepackage{xspace}
\usepackage{cancel}
\usepackage{mathtools}
\usepackage{amsthm}
\makeatletter
\@ifundefined{theorem}{\newtheorem{theorem}{Theorem}}{}
\@ifundefined{lemma}{\newtheorem{lemma}[theorem]{Lemma}}{}
\@ifundefined{proposition}{\newtheorem{proposition}[theorem]{Proposition}}{}
\makeatother
\title[Density functions for filtrations of graded ideals]{Density functions for filtrations of graded ideals}
\author{Suprajo Das, Hoang Le Truong}
\date{Source: arXiv:2609.18335v1 (CC BY 4.0)}
\begin{document}
\maketitle

\noindent\emph{Source and licence.} Density functions for filtrations of graded ideals. Version arXiv:2609.18335v1 (CC BY 4.0), distributed under the Creative Commons Attribution 4.0 International licence, as recorded in the arXiv licence metadata for this exact version and shown on the abstract page https://arxiv.org/abs/2609.18335v1. Full rights notice: licenses/arxiv-2609.18335-CC-BY-4.0.txt.\par\smallskip

\noindent\textit{Editorial scope.} Counted unit: the Proposition whose source label is \texttt{welldefined} in the article (source0.tex lines 319--326, source label \texttt{welldefined}), delivered together with the article's own complete original proof of it (source0.tex lines 327--397, one \texttt{proof} environment of 71 lines). No mathematics is written, repaired or re-derived: statement and proof are the article's own text line for line. Required context retained uncounted: lem\_converge (lines 296--298). Every definition, display and label of those lines is retained unchanged. Omitted from this package and not redistributed: 1--295, 299--318, 398--1149. Nothing inside a retained excerpt is omitted or altered: every retained body is delivered exactly as the article's lines are written, so no omission receipt of any kind accompanies this package. Citations: the retained excerpts carry no citation command at all, so no bibliography entry of the article is transcribed and no reference is redistributed. Reference adaptation: none was needed and none was made; every reference inside the delivered unit resolves inside it, so no unresolved reference or citation is produced. Printed slips kept literal: no printed word, symbol, hypothesis or number of the counted statement or of its proof was altered. Layout and class: the article's own class is \texttt{amsart} with packages including enumerate, amsmath, amsthm, verbatim, amssymb, amsfonts, amscd, graphicx, mathtools, graphics, hyperref, mathrsfs, geometry, inputenc; the wrapper is the lane's pinned \texttt{amsart} editorial wrapper at 10pt on A4, which restates the theorem-like environments the retained text needs and adds the article's own macro definitions that the retained text needs (none), with extra wrapper packages (bm, bbm, dsfont, xspace, cancel, mathtools). The delivered numbering is the unit's own. Assets: no retained span contains a figure environment, a graphics-inclusion command, a PDF-page inclusion, a table or a TikZ picture; no image file is copied into this package, the package declares no asset dependency and no image file is redistributed with it. Licence: version arXiv:2609.18335v1 of this article is distributed under the Creative Commons Attribution 4.0 International licence, as recorded in the arXiv licence metadata for this exact version and shown on the abstract page https://arxiv.org/abs/2609.18335v1. A full rights notice is delivered at licenses/arxiv-2609.18335-CC-BY-4.0.txt. Characters: every retained excerpt is pure ASCII, so the delivered engine, XeLaTeX with its default Latin Modern text font, reports no missing character.
\begin{lemma}\label{lem_converge}
Let $p,q>0$ be integers such that $\tfrac{p}{q}>\alpha_{\mathbb{I}}$. Then there exists an integer $n_0>0$ such that ${\left(I_{qn}\right)}_{pn-1}\neq 0$ and ${\left(I_{qn+1}\right)}_{pn}\neq 0$ for all integers $n\geq n_0$.
\end{lemma}

\begin{proposition}\label{welldefined}
Let $p,q>0$ be integers such that $\tfrac{p}{q}>\alpha_{\mathbb{I}}$. Then
\begin{align*}
G\big(\Gamma^{\mathbb{I}}_{(p,q)}\big) \cap \left(\mathbb{Z}^d\times \{0\}\right) &= G\big(\Gamma^{R_{\bullet}}_{(1,1)}\big) \cap \left(\mathbb{Z}^d\times \{0\}\right),\;\;\text{and}\\
L\big(\Gamma^{\mathbb{I}}_{(p,q)}\big) \cap \left(\mathbb{Z}^d\times \{0\}\right) &= L\big(\Gamma^{R_{\bullet}}_{(1,1)}\big) \cap \left(\mathbb{Z}^d\times \{0\}\right).
\end{align*}
In particular, $\mathrm{ind}\big(\Gamma^{\mathbb{I}}_{(p,q)}\big)= \mathrm{ind}\big(\Gamma^{{R}_{\bullet}}_{(1,1)}\big)$, where ${R}_{\bullet}$ denotes the constant filtration $\{R\}_{n\in\mathbb{N}}$. Moreover, for any integer $n_0>0$, we have $\frac{1}{q}\cdot\Delta\big(\Gamma^{\mathbb{I}}_{(p,q)}\big) = \frac{1}{qn_0}\cdot\Delta\big(\Gamma^{\mathbb{I}}_{(pn_0,qn_0)}\big)$.
\end{proposition}

\begin{proof} We shall breakdown our proof into multiple claims.

\vspace{0.2cm}
\noindent
\emph{Claim $1$}. For any integer $n_0>0$, we have
\begin{equation*}
L\big(\Gamma^{\mathbb{I}}_{(p,q)}\big) = L\big(\Gamma^{\mathbb{I}}_{(pn_0,qn_0)}\big),\;\; \mathrm{Con}\big(\Gamma^{\mathbb{I}}_{(p,q)}\big) = \mathrm{Con}\big(\Gamma^{\mathbb{I}}_{(pn_0,qn_0)}\big), \;\; \text{and} \;\; \frac{1}{q}\cdot\Delta\big(\Gamma^{\mathbb{I}}_{(p,q)}\big) = \frac{1}{qn_0}\cdot\Delta\big(\Gamma^{\mathbb{I}}_{(pn_0,qn_0)}\big).
\end{equation*}

For each $n_0>0$, there are inclusions $$n_0 \cdot \Gamma^{\mathbb{I}}_{(p,q)} \subseteq \Gamma^{\mathbb{I}}_{(pn_0,qn_0)} \subseteq \Gamma^{\mathbb{I}}_{(p,q)}.$$ These imply that $L\big(\Gamma^{\mathbb{I}}_{(p,q)}\big) = L\big(\Gamma^{\mathbb{I}}_{(pn_0,qn_0)}\big)$ and $\mathrm{Con}\big(\Gamma^{\mathbb{I}}_{(p,q)}\big) = \mathrm{Con}\big(\Gamma^{\mathbb{I}}_{(pn_0,qn_0)}\big)$. Intersecting with the hyperplane $\mathbb{R}^d\times \{1\}$ yields
\begin{multline*}
 \frac{1}{q}\cdot\Delta\big(\Gamma^{\mathbb{I}}_{(p,q)}\big) = \frac{1}{q}\cdot\left(\mathrm{Con}\big(\Gamma^{\mathbb{I}}_{(p,q)}\big) \cap \left(\mathbb{R}^d\times \{q\}\right)\right) = \mathrm{Con}\big(\Gamma^{\mathbb{I}}_{(p,q)}\big) \cap \left(\mathbb{R}^d\times \{1\}\right)\\ =\mathrm{Con}\big(\Gamma^{\mathbb{I}}_{(pn_0,qn_0)}\big) \cap \left(\mathbb{R}^d\times \{1\}\right) = \frac{1}{qn_0}\cdot\left(\mathrm{Con}\big(\Gamma^{\mathbb{I}}_{(pn_0,qn_0)}\big) \cap \left(\mathbb{R}^d\times \{qn_0\}\right)\right) = \frac{1}{qn_0}\cdot\Delta\big(\Gamma^{\mathbb{I}}_{(pn_0,qn_0)}\big).
\end{multline*}

\vspace{0.2cm}
\noindent
\emph{Claim $2$}. For every integer $n_0>0$, we have
\begin{equation*}
G\big(\Gamma^{\mathbb{I}}_{(p,q)}\big) \cap \left(\mathbb{Z}^d\times \{0\}\right) = G\big(\Gamma^{\mathbb{I}}_{(pn_0,qn_0)}\big) \cap \left(\mathbb{Z}^d\times \{0\}\right).
\end{equation*}

Clearly, $G\big(\Gamma^{\mathbb{I}}_{(pn_0,qn_0)}\big) \cap \left(\mathbb{Z}^d\times \{0\}\right) \subseteq G\big(\Gamma^{\mathbb{I}}_{(p,q)}\big) \cap \left(\mathbb{Z}^d\times \{0\}\right)$.
Let $(\mathbf{u},0) \in G\big(\Gamma^{\mathbb{I}}_{(p,q)}\big) \cap \left(\mathbb{Z}^d\times \{0\}\right)$. Then there exist $\mathbf{a}, \mathbf{b}\in\mathbb{N}^d$, and $s\in\mathbb{N}$ such that $(\mathbf{u},0) = (\mathbf{a}, qs) - (\mathbf{b}, qs)$ with $(\mathbf{a}, qs), (\mathbf{b}, qs) \in \Gamma^{\mathbb{I}}_{(p,q)}$. As ${(I_q)}_p\neq 0$, there exist $\mathbf{c}\in\mathbb{N}^d$ with $(\mathbf{c},q)\in \Gamma^{\mathbb{I}}_{(p,q)}$. Write $s = n_0k + t$, where $0\leq t\leq n_0-1$. So,
\begin{align*}
 (\mathbf{u},0) &= (\mathbf{a}, qs) - (\mathbf{b}, qs)\\
 &= (\mathbf{a}, qs) + (n_0-t)\cdot (\mathbf{c},q) - (\mathbf{b}, qs) - (n_0-t)\cdot (\mathbf{c},q)\\
 &= (\mathbf{a} + (n_0-t)\mathbf{c}, (k+1)qn_0) - (\mathbf{b} + (n_0-t)\mathbf{c}, (k+1)qn_0),
\end{align*}
where $(\mathbf{a} + (n_0-t)\mathbf{c}, (k+1)qn_0), (\mathbf{b} + (n_0-t)\mathbf{c}, (k+1)qn_0) \in \Gamma^{\mathbb{I}}_{(pn_0,qn_0)}$. In other words, $(\mathbf{u},0) \in G\big(\Gamma^{\mathbb{I}}_{(pn_0,qn_0)}\big) \cap \left(\mathbb{Z}^d\times \{0\}\right)$. This completes the proof of Claim $2$.

\vspace{0.2cm}
\noindent
\emph{Claim $3$}. Let $p^{\prime},q^{\prime}>0$ be integers such that $\frac{p^{\prime}}{q^{\prime}}>\frac{p}{q}$. Then
\begin{align}
G\big(\Gamma^{\mathbb{I}}_{(p,q)}\big) \cap \left(\mathbb{Z}^d\times \{0\}\right) &= G\big(\Gamma^{\mathbb{I}}_{(p^{\prime},q^{\prime})}\big) \cap \left(\mathbb{Z}^d\times \{0\}\right),\;\;\text{and}\label{equal1}\\
L\big(\Gamma^{\mathbb{I}}_{(p,q)}\big) \cap \left(\mathbb{Z}^d\times \{0\}\right) &= L\big(\Gamma^{\mathbb{I}}_{(p^{\prime},q^{\prime})}\big) \cap \left(\mathbb{Z}^d\times \{0\}\right).\label{equal2}
\end{align}

Choose $n_0, n_0^{\prime}>0$ with $qn_0 = q^{\prime}n_0^{\prime}$. By Claim $1$ and Claim $2$, we may replace $(p,q)$ and $(p^{\prime},q^{\prime})$ by $(pn_0,qn_0)$ and $(p^{\prime}n^{\prime}_0,q^{\prime}n^{\prime}_0)$ respectively, and by Lemma \ref{lem_converge}, we may further assume that ${\left(I_{q}\right)}_{p-1} \neq 0$. Thus, it suffices to assume from the beginning that $q=q^{\prime}$, and ${\left(I_{q}\right)}_{p-1} \neq 0$.

Let $z\in R_1$ and $f\in {\left(I_{q}\right)}_{p-1}$ be two nonzero elements. There exist unique vectors $\mathbf{v}_z, \mathbf{v}_f \in \mathbb{N}^d$ such that $\nu(z) = \langle \mathbf{v}_z, \boldsymbol{\xi}\rangle$, and $\nu(f) = \langle \mathbf{v}_f, \boldsymbol{\xi}\rangle$, where $\boldsymbol{\xi}=(\xi_1,\ldots,\xi_d)\in \mathbb{R}^d$. Put $c=p^{\prime}-p$. For all $n,n^{\prime}>0$, there are inclusions $$z^{cn}\cdot {\left(I_{qn}\right)}_{pn} \subseteq {\left(I_{qn}\right)}_{p^{\prime}n}, \quad \text{and} \quad f^{cn^{\prime}}\cdot {\left(I_{qn^{\prime}}\right)}_{p^{\prime}n^{\prime}} \subseteq {\left(I_{(1+c)qn^{\prime}}\right)}_{(1+c)pn^{\prime}}.$$ In other words, if $(\mathbf{u},qn)\in \Gamma^{\mathbb{I}}_{(p,q)}$ then $(\mathbf{u}+cn\mathbf{v}_z,qn)\in \Gamma^{\mathbb{I}}_{(p^{\prime},q)}$. Similarly, if $(\mathbf{u}^{\prime},qn^{\prime})\in \Gamma^{\mathbb{I}}_{(p^{\prime},q)}$ then $(\mathbf{u}^{\prime}+cn^{\prime}\mathbf{v}_f,(1+c)qn^{\prime})\in \Gamma^{\mathbb{I}}_{\left((1+c)p,(1+c)q\right)}$.

To prove \eqref{equal1}, let $\left(\mathbf{u},0\right) \in G\big(\Gamma^{\mathbb{I}}_{(p,q)}\big) \cap \left(\mathbb{Z}^d\times \{0\}\right)$. Then there exist $\mathbf{a},\mathbf{b}\in \mathbb{N}^d$, and $n\in\mathbb{N}$ such that $\left(\mathbf{u},0\right) = (\mathbf{a},qn) - (\mathbf{b},qn)$, with $(\mathbf{a},qn), (\mathbf{b},qn) \in \Gamma^{\mathbb{I}}_{(p,q)}$. Since we can write $(\mathbf{a},qn) - (\mathbf{b},qn) = (\mathbf{a}+cn\mathbf{v}_z,qn) - (\mathbf{b}+cn\mathbf{v}_z,qn)$, it follows that $\left(\mathbf{u},0\right) \in G\big(\Gamma^{\mathbb{I}}_{(p^{\prime},q)}\big) \cap \left(\mathbb{Z}^d\times \{0\}\right)$. Again, let $\left(\mathbf{u}^{\prime},0\right) \in G\big(\Gamma^{\mathbb{I}}_{(p^{\prime},q)}\big) \cap \left(\mathbb{Z}^d\times \{0\}\right)$. Then we can write $\left(\mathbf{u}^{\prime},0\right) = (\mathbf{a}^{\prime},qn^{\prime}) - (\mathbf{b}^{\prime},qn^{\prime})$, for some $(\mathbf{a}^{\prime},qn^{\prime}), (\mathbf{b}^{\prime},qn^{\prime}) \in \Gamma^{\mathbb{I}}_{(p^{\prime},q)}$. Since $(\mathbf{a}^{\prime},qn^{\prime}) - (\mathbf{b}^{\prime},qn^{\prime}) = (\mathbf{a}^{\prime}+cn^{\prime}\mathbf{v}_f,qn^{\prime}) - (\mathbf{b}^{\prime}+cn^{\prime}\mathbf{v}_f,qn^{\prime})$, it follows that $\left(\mathbf{u}^{\prime},0\right) \in G\big(\Gamma^{\mathbb{I}}_{((1+c)p,(1+c)q)}\big)\cap \left(\mathbb{Z}^d\times \{0\}\right)$. So, $$G\big(\Gamma^{\mathbb{I}}_{(p,q)}\big) \cap \left(\mathbb{Z}^d\times \{0\}\right) \subseteq G\big(\Gamma^{\mathbb{I}}_{(p^{\prime},q)}\big) \cap \left(\mathbb{Z}^d\times \{0\}\right) \subseteq G\big(\Gamma^{\mathbb{I}}_{((1+c)p,(1+c)q)}\big) \cap \left(\mathbb{Z}^d\times \{0\}\right).$$ Clearly, $\Gamma^{\mathbb{I}}_{((1+c)p,(1+c)q)} \subseteq \Gamma^{\mathbb{I}}_{(p,q)}$ and this proves \eqref{equal1}. One also proves \eqref{equal2} by similar methods.

\vspace{0.2cm}
\noindent
\emph{Claim $4$}. We have
\begin{align*}
G\big(\Gamma^{\mathbb{I}}_{(p,q)}\big) \cap \left(\mathbb{Z}^d\times \{0\}\right) &= G\big(\Gamma^{R_{\bullet}}_{(1,1)}\big) \cap \left(\mathbb{Z}^d\times \{0\}\right),\;\;\text{and}\\
L\big(\Gamma^{\mathbb{I}}_{(p,q)}\big) \cap \left(\mathbb{Z}^d\times \{0\}\right) &= L\big(\Gamma^{R_{\bullet}}_{(1,1)}\big) \cap \left(\mathbb{Z}^d\times \{0\}\right).
\end{align*}

By Claim $3$, we may reduce to $q=1$ and $(I_1)_{p-1}\neq 0$. Pick a nonzero element $g \in {\left(I_{1}\right)}_{p-1}$. There is a unique vector $\mathbf{v}_g \in \mathbb{N}^d$ such that $\nu(g) = \langle \mathbf{v}_g, \boldsymbol{\xi}\rangle$. Let $g_{\bullet}$ denote the filtration $\left\{(g^n)\right\}_{n\in\mathbb{N}}$. Then
\begin{align*}
\Gamma^{g_{\bullet}}_{(p,1)} &= \left\{(\mathbf{v}_f,n)\in\mathbb{N}^{d+1} \mid \exists f \in {(g^{n})}_{pn}\;\text{with}\;\nu(f) = \langle \mathbf{v}_f, \boldsymbol{\xi} \rangle\right\}\\
  &= \left\{(n\mathbf{v}_g + \mathbf{v}_{f^{\prime}},n)\in\mathbb{N}^{d+1} \mid \exists f^{\prime} \in {R}_{n}\;\text{with}\;\nu(f^{\prime}) = \langle \mathbf{v}_{f^{\prime}}, \boldsymbol{\xi} \rangle\right\},
\end{align*}
which shows that
\begin{align*}
 G\big(\Gamma^{g_{\bullet}}_{(p,1)}\big)\cap \left(\mathbb{Z}^d\times \{0\}\right) &= G\big(\Gamma^{R_{\bullet}}_{(1,1)}\big)\cap \left(\mathbb{Z}^d\times \{0\}\right),\;\text{and}\\
  L\big(\Gamma^{g_{\bullet}}_{(p,1)}\big)\cap \left(\mathbb{Z}^d\times \{0\}\right) &= L\big(\Gamma^{R_{\bullet}}_{(1,1)}\big)\cap \left(\mathbb{Z}^d\times \{0\}\right).
 \end{align*}
For each $n$, we have $(g^n)_{pn} \subseteq (I_n)_{pn} \subseteq R_{pn}$. Now using the above equalities, one gets
\begin{align*}
 G\big(\Gamma^{R_{\bullet}}_{(1,1)}\big)\cap \left(\mathbb{Z}^d\times \{0\}\right) &\subseteq G\big(\Gamma^{\mathbb{I}}_{(p,1)}\big)\cap \left(\mathbb{Z}^d\times \{0\}\right) \subseteq G\big(\Gamma^{R_{\bullet}}_{(p,1)}\big)\cap \left(\mathbb{Z}^d\times \{0\}\right),\;\text{and}\\
 L\big(\Gamma^{R_{\bullet}}_{(1,1)}\big)\cap \left(\mathbb{Z}^d\times \{0\}\right) &\subseteq L\big(\Gamma^{\mathbb{I}}_{(p,1)}\big)\cap \left(\mathbb{Z}^d\times \{0\}\right) \subseteq L\big(\Gamma^{R_{\bullet}}_{(p,1)}\big)\cap \left(\mathbb{Z}^d\times \{0\}\right).
\end{align*}
The outer terms are equal from Claim $3$, hence all coincide. Thus Claim $4$ is established.

The equality $\mathrm{ind}\big(\Gamma^{\mathbb{I}}_{(p,q)}\big) = \mathrm{ind}\big(\Gamma^{R_{\bullet}}_{(1,1)}\big)$ is a direct consequence of Claim $4$.
\end{proof}

\end{document}
